\section*{1. Adversarial Alien Obfuscation}

\subsection*{Problem \& Approach}
The goal is to determine if renaming alien operators to semantic strings (e.g., \texttt{combine}) or designed character strings affects the model's ability to reason. We evaluate robustness by replacing the original alien operator names with four obfuscation schemes:
\begin{itemize}
    \item \textbf{Semantic}: Similar idea English words (e.g., \texttt{combine}).
    \item \textbf{Misleading}: Unrelated English words (e.g., \texttt{paint}).
    \item \textbf{Gibberish}: Random strings (e.g., \texttt{qahft}).
    \item \textbf{Unicode}: Non-ASCII symbols (e.g., $\Delta$).
\end{itemize}

\begin{table}[h]
\centering
\begin{tabular}{l c c c c}
\toprule
\textbf{Scheme} & \textbf{Overall Acc.} & \textbf{Level 1} & \textbf{Level 2} & \textbf{Level 3} \\
\midrule
Original & 11.3\% & 30.0\% & 2.0\% & 2.0\% \\
Gibberish & 10.0\% & 28.0\% & 0.0\% & 2.0\% \\
Unicode & 6.0\% & 16.0\% & 0.0\% & 2.0\% \\
Semantic & 0.0\% & 0.0\% & 0.0\% & 0.0\% \\
Misleading & 0.0\% & 0.0\% & 0.0\% & 0.0\% \\
\bottomrule
\end{tabular}
\caption{Model accuracy under various obfuscation schemes.}
\label{tab:obfuscation_results}
\end{table}

In the Semantic and Misleading schemes, the attack worked.
With meaningless symbols (Original, Gibberish, Unicode), the model could have still treated the prompt as an logic puzzle, solved it and followed the formatting to output exactly \texttt{Final Answer: <number>}, but with real English words it likely triggered conversational priors from training data. The model outputs natural sentences (e.g., \texttt{The final answer is 23.}), which even though correct couldn't be extracted by regex matching and thus marked as failed. 
